\documentclass[addpoints]{exam}

\usepackage{amsmath}
\usepackage{amssymb}
\usepackage{amsfonts}
\usepackage{mathtools}
\usepackage{enumitem}
\usepackage{tikz}

\DeclarePairedDelimiter{\Norm}\lVert\rVert
\newcommand{\R}{\mathbb{R}}
\newcommand{\tr}{\intercal}
\newcommand{\im}{\mathop{\mathrm{im}}}

%EDIT THESE WHEN NEW CLASS
\newcommand{\themath}{19620}
\newcommand{\thesec}{46}

\pagestyle{headandfoot}
\runningheader{MATH \themath\ Section \thesec}{Final Prototype}{Duarte Maia}
\runningheadrule
\footer{}{\ifincomplete{Exercise continues on next page.}{}}{}


\begin{document}

\begin{center}
\Large Final Prototype

MATH \themath - Section \thesec

\smallskip

\large Section Instructor: Duarte Maia

\smallskip

March 10, 2026
\end{center}

\begin{center}
\fbox{
\parbox{6in}{
\begin{itemize}
\item You have 2 hours to answer the questions below on the provided examination book.

\item Write your name and UCID in the cover page of the examination book.

\item You may ask for additional paper. If you do, sign both exam books and number them.

\item Unless otherwise specified, you do not need to justify your answers. However, you must still show your work. You may lose points if I deem you to have performed a significant leap in logic.

\item This exam is graded out of 100 points. [Note: The prototype questions add up to 130 points, because only three questions from Part I will be on the exam.]

\item Errors in reasoning will be penalized more severely than errors in calculation. If you make a mistake in your calculations but interpret the results correctly, you will get most of the points for the question, unless your calculation error bypasses a significant amount of reasoning you would otherwise have to do.
\end{itemize}
}
}
\end{center}


\noindent Useful Formulas:
\begin{itemize}
\item $\det(AB) = \det A \det B$, $\det(A^{-1}) = 1/\det A$, $\det A^\tr = \det A$.
\item Laplace Expansion: $\det A = \sum_i (-1)^{i+j} a_{ij} \det A_{ij}$.
\item Quadratic formula: $ax^2 + bx + c = 0$ solved by $x = \frac{-b\pm\sqrt{b^2 -4ac}}{2a}$.
\item Projection on line spanned by unit vector $\vec u$: $\vec u \vec u^\tr$.
\item Angle between $\vec v$ and $\vec w$: $\theta = \arccos\left(\frac{\vec v \cdot \vec w}{\Norm{\vec v} \Norm{\vec w}}\right)$.
\end{itemize}

%\hrule

\qformat{\textbf{Exercise \thequestion} \quad (Total: \totalpoints \ points)\hfill}

\newpage 
\begin{questions}
\fullwidth{\Large\bf Part I -- General Linear Algebra}
\fullwidth{\textit{Note: In the real exam, only three of the following question types would appear.}}
\question Let $A$ be a $5\times 3$ matrix, meaning it has 5 rows and 3 columns.

\begin{parts}
\part[3] What are the possible values of $\dim(\ker A)$? You do not need to justify your answer.
\part[3] What are the possible values of $\dim(\im A)$? You do not need to justify your answer.
\part[4] Suppose that the subset of $\R^5$ defined by $A^\tr \vec x = \vec 0$ is a two-dimensional plane. True or false: ``The equation $A^\tr \vec x = \vec b$ has a solution, no matter the choice of $\vec b$.'' Justify your answer.
\end{parts}

\question\label{q:v} Let $V$ be the subspace of $\R^5$ spanned by:
\[\vec v_1 = \begin{bmatrix} 1 \\ 2 \\ 0 \\ 1 \\ 2 \end{bmatrix},\quad
\vec v_2 = \begin{bmatrix} -1 \\ -2 \\ 2 \\ 0 \\ -2 \end{bmatrix},\quad
\vec v_3 = \begin{bmatrix} 0 \\ 0 \\ 2 \\ 1 \\ 0 \end{bmatrix}, \text{ and } \quad
\vec v_4 = \begin{bmatrix} 1 \\ 2 \\ 2 \\ 3 \\ 4 \end{bmatrix}.\]
\begin{parts}
\part[5] Calculate the dimension of $V$.
\part[5] Find a spanning set of $V$ with five distinct elements.
\end{parts}

\question Consider the matrix
\[A = \begin{bmatrix}
3 & -2 & 6 \\
-1 & 1 & -2 \\
-1 & 1 & -3
\end{bmatrix}.\]
\begin{parts}
\part[7] Calculate $A^{-1}$.
\part[3] \textit{Using the result from the previous part}, solve the equation $A\vec x = (1,1,1)$.
\end{parts}

\question Consider the four vectors
\[\vec v_1 = \begin{bmatrix} 1 \\ 2 \\ 0 \\ 1 \\ 2 \end{bmatrix},\quad
\vec v_2 = \begin{bmatrix} -1 \\ -2 \\ 2 \\ 0 \\ -2 \end{bmatrix},\quad
\vec v_3 = \begin{bmatrix} 0 \\ 0 \\ 2 \\ 1 \\ 0 \end{bmatrix}, \text{ and } \quad
\vec v_4 = \begin{bmatrix} 1 \\ 2 \\ 2 \\ 3 \\ 4 \end{bmatrix}.\]
\textit{These are the same vectors from Exercise \ref{q:v}. This means that some of your work can possibly be recycled.}
\begin{parts}
\part[5] Is $\vec v_4$ in $\mathrm{span}(\vec v_1, \vec v_2, \vec v_3)$?
\part[5] Is the list $(\vec v_1, \vec v_2, \vec v_3, \vec v_4)$ linearly dependent or independent?
\end{parts}

\question Sketch the following subspaces of $\R^2$.
\begin{parts}
\part[2] The subspace spanned by $(999,1)$, $(1,777)$, and $(123,456)$,
\part[4] The kernel of $\left[\begin{smallmatrix} 1 & 2 \\ 3 & 4 \end{smallmatrix}\right]$,
\part[4] The image of $\left[\begin{smallmatrix} 1 & 1 \\ 1 & 1 \end{smallmatrix}\right]$.
\end{parts}

\question[10] Consider the following matrices:
\[A = \begin{bmatrix} 1 & -1 \\ -1 & 1 \end{bmatrix}, \quad
B = \begin{bmatrix} 2 & -2 \\ 1 & -1 \\ 3 & -3 \end{bmatrix},\text{ and }\quad
C = \begin{bmatrix} 1 & 1 \\ 1 & 1 \end{bmatrix}.\]
Each of these three matrices corresponds to two different subspaces: The kernel and the image. Which of these six subspaces are the same, and which are distinct?

\fullwidth{\Large\bf Part II -- Geometry}

\question[5] Consider the linear transformation
\[T(x,y,z)=(x+2y,2z+x,3z+x+y).\]
A cube $Q$ of side length $2$ is distorted by applying $T$, as to obtain the shape $TQ$. What is the volume of $TQ$?

\question Let $A$ be a $4\times 4$ matrix whose determinant is equal to $-2$. Let $S$ be an invertible $4\times 4$ matrix of unknown determinant $\det S$.
\begin{parts}
\part[2] Calculate $\det(SAS^{-1})$. Justify your answer.
\part[3] Calculate $\det(-A)$. Justify your answer.
\end{parts}
 
\question[5] Find the angle between two edges of a regular tetrahedron (that is, a triangular pyramid whose edges all have the same length).

\textit{Hint: The points $(1,1,-1)$, $(1,-1,1)$, $(-1,1,1)$, and $(-1,-1,-1)$ form the vertices of a regular tetrahedron. You may use this fact without justification.}

\question[5] Let $V$ be the subspace of $\R^4$ defined by the equations:
\[\left\{
\arraycolsep=1.4pt
\begin{array}{rrrrrrrrrr}
x&+&y&+&z&+&w&=&0,\\
x&-&y&+&2z&+&w&=&0.
\end{array}
\right.\]
Show that every vector in $V$ is orthogonal to the vector $(2,0,3,2)$.


\fullwidth{\Large\bf Part III -- Diagonalization}

\question[20] Diagonalize the matrix
\[A =
\begin{bmatrix}
1 & 0 & 0 \\
-2 & -1 & 2 \\
0 & 0 & 1
\end{bmatrix}.\]

\question[20] Two antagonistic countries, $X$ and $Y$, adjust their defense budget monthly in response to the other. Let $x(t)$ and $y(t)$ be their defense budget, in millions of dollars, $t$ months after some fixed date.

Country $X$ is peaceful by default, but is scared of $Y$, and will therefore adjust its defense budget downward unless $Y$ has is looking scary. Let's say $x(t+1)=0.6 x(t) + 0.8 y(t)$.

Country $Y$ has ambitions of world domination, but these ambitions are tempered if country $X$ is looking scary. Let's say $y(t+1)=1.2 y(t) - 0.1 x(t)$. If this calculation ever produces a negative number, let's say that country $Y$ has yielded to the overwhelming power of $X$ and agreed to a disarmament treaty.

\begin{parts}
\part[4] Let $\vec m(t) = (x(t),y(t))$. Find a matrix $A$ such that $\vec m(t+1)= A \vec m(t)$.

\part[8] The matrix $A$ you found in the previous exercise can be diagonalized as $A = SDS^{-1}$, where
\[
S = \begin{bmatrix}
2 & 4\\
1 & 1
\end{bmatrix},
\quad
D = \begin{bmatrix}
1 & 0\\
0 & 0.8
\end{bmatrix}
,\text{ and} \quad
S^{-1} = \begin{bmatrix}
-\tfrac12 & 2 \\
\tfrac12 & -1
\end{bmatrix}.
\]

If countries $X$ and $Y$ hold fast in their policy, what is the fate of the world? Indefinite escalation? Disarmament? An uneasy peace?

There are multiple cases, depending on the initial defense budgets $x(0)$ and $y(0)$. Describe them, including edge cases.

\part[8] In an alternate universe, country $X$ is slightly less reactive, following instead the rule $x(t+1) = 0.6 x(t) + 0.5 y(t)$. In this case, the matrix of the system diagonalizes as $S' D' (S')^{-1}$, with
\[
S' = \begin{bmatrix}
1 & 5\\
1 & 1
\end{bmatrix},
\quad
D' = \begin{bmatrix}
1.1 & 0\\
0 & 0.6
\end{bmatrix}
,\text{ and} \quad
(S')^{-1} = \frac14\begin{bmatrix}
-1 & 5 \\
1 & -1
\end{bmatrix}.
\]

As before, determine the fate of the world, as a function of the initial defense budgets.
\end{parts}

\question[10] Some of the following statements are true, and some are false. Justify the true statement(s) and find counter-example(s) for the false statement(s).
\begin{enumerate}[label={(\roman*)}]
\item If $\vec v$ is an eigenvector of $A^3$ with eigenvalue $\lambda$, then $\vec v$ is also an eigenvector of $A$ with eigenvalue $\sqrt[3]\lambda$.
\item If $\vec v$ is an eigenvector of $A$, then $\vec v$ is also an eigenvector of $A^\tr$.
\item If $A$ is a diagonalizable matrix all of whose eigenvalues are in the interval $(-1,1)$, then $\lim_{n\to\infty} A^n$ is the zero matrix.
\end{enumerate}

\end{questions}
\end{document}